% Options for packages loaded elsewhere
\PassOptionsToPackage{unicode}{hyperref}
\PassOptionsToPackage{hyphens}{url}
\PassOptionsToPackage{dvipsnames,svgnames,x11names}{xcolor}
\documentclass[
  11pt,
]{article}
\usepackage{xcolor}
\usepackage[a4paper,margin=0.8in]{geometry}
\usepackage{amsmath,amssymb}
\setcounter{secnumdepth}{5}
\usepackage{iftex}
\ifPDFTeX
  \usepackage[T1]{fontenc}
  \usepackage[utf8]{inputenc}
  \usepackage{textcomp} % provide euro and other symbols
\else % if luatex or xetex
  \usepackage{unicode-math} % this also loads fontspec
  \defaultfontfeatures{Scale=MatchLowercase}
  \defaultfontfeatures[\rmfamily]{Ligatures=TeX,Scale=1}
\fi
\usepackage{lmodern}
\ifPDFTeX\else
  % xetex/luatex font selection
  \setmainfont[]{Times New Roman}
  \setsansfont[]{Calibri}
  \setmonofont[]{Consolas}
\fi
% Use upquote if available, for straight quotes in verbatim environments
\IfFileExists{upquote.sty}{\usepackage{upquote}}{}
\IfFileExists{microtype.sty}{% use microtype if available
  \usepackage[]{microtype}
  \UseMicrotypeSet[protrusion]{basicmath} % disable protrusion for tt fonts
}{}
\makeatletter
\@ifundefined{KOMAClassName}{% if non-KOMA class
  \IfFileExists{parskip.sty}{%
    \usepackage{parskip}
  }{% else
    \setlength{\parindent}{0pt}
    \setlength{\parskip}{6pt plus 2pt minus 1pt}}
}{% if KOMA class
  \KOMAoptions{parskip=half}}
\makeatother
\usepackage{longtable,booktabs,array}
\newcounter{none} % for unnumbered tables
\usepackage{calc} % for calculating minipage widths
% Correct order of tables after \paragraph or \subparagraph
\usepackage{etoolbox}
\makeatletter
\patchcmd\longtable{\par}{\if@noskipsec\mbox{}\fi\par}{}{}
\makeatother
% Allow footnotes in longtable head/foot
\IfFileExists{footnotehyper.sty}{\usepackage{footnotehyper}}{\usepackage{footnote}}
\makesavenoteenv{longtable}
\usepackage{graphicx}
\makeatletter
\newsavebox\pandoc@box
\newcommand*\pandocbounded[1]{% scales image to fit in text height/width
  \sbox\pandoc@box{#1}%
  \Gscale@div\@tempa{\textheight}{\dimexpr\ht\pandoc@box+\dp\pandoc@box\relax}%
  \Gscale@div\@tempb{\linewidth}{\wd\pandoc@box}%
  \ifdim\@tempb\p@<\@tempa\p@\let\@tempa\@tempb\fi% select the smaller of both
  \ifdim\@tempa\p@<\p@\scalebox{\@tempa}{\usebox\pandoc@box}%
  \else\usebox{\pandoc@box}%
  \fi%
}
% Set default figure placement to htbp
\def\fps@figure{htbp}
\makeatother
\setlength{\emergencystretch}{3em} % prevent overfull lines
\providecommand{\tightlist}{%
  \setlength{\itemsep}{0pt}\setlength{\parskip}{0pt}}

% ── Better tables ──
\usepackage{booktabs}
\renewcommand{\arraystretch}{1.3}

% ── Figure handling ──
\usepackage{float}

%  Force max image size 
\usepackage{graphicx}
\setkeys{Gin}{width=0.95\linewidth,keepaspectratio}


% ── Colors ──
\usepackage{xcolor}
\definecolor{darkblue}{RGB}{0,51,102}

% ── Section heading formatting ──
\usepackage{titlesec}
\titleformat{\section}{\Large\bfseries\sffamily\color{darkblue}}{\thesection}{1em}{}
\titleformat{\subsection}{\large\bfseries\sffamily\color{darkblue}}{\thesubsection}{1em}{}
\titleformat{\subsubsection}{\normalsize\bfseries\sffamily\color{darkblue}}{\thesubsubsection}{1em}{}
\titlespacing*{\section}{0pt}{18pt}{8pt}
\titlespacing*{\subsection}{0pt}{14pt}{6pt}
\titlespacing*{\subsubsection}{0pt}{10pt}{4pt}

% ── Paragraph spacing ──
\setlength{\parskip}{6pt plus 2pt minus 1pt}
\setlength{\parindent}{0pt}

% ── Header/footer ──
\usepackage{fancyhdr}
\pagestyle{fancy}
\fancyhf{}
\fancyhead[L]{\small\sffamily\color{gray}Solar PV Forecasting -- Journal Paper}
\fancyhead[R]{\small\sffamily\color{gray}\thepage}
\fancyfoot[C]{}
\renewcommand{\headrulewidth}{0.4pt}

% ── Prevent widows/orphans ──
\widowpenalty=10000
\clubpenalty=10000
\usepackage{bookmark}
\IfFileExists{xurl.sty}{\usepackage{xurl}}{} % add URL line breaks if available
\urlstyle{same}
\hypersetup{
  colorlinks=true,
  linkcolor={darkblue},
  filecolor={Maroon},
  citecolor={Blue},
  urlcolor={darkblue},
  pdfcreator={LaTeX via pandoc}}

\author{}
\date{}

\begin{document}

{
\setcounter{tocdepth}{3}
\tableofcontents
}
\section{Forecasting the Output of Solar Photovoltaic Plants Using
Statistical Analysis and Machine Learning: An Hourly Resolution
Study}\label{forecasting-the-output-of-solar-photovoltaic-plants-using-statistical-analysis-and-machine-learning-an-hourly-resolution-study}

\textbf{Authors:} Abdul Saleem Shaik, {[}Co-Author Name{]}, {[}Co-Author
Name{]} \textbf{Affiliation:} {[}Department of Electrical/Energy
Engineering, University Name, Country{]} \textbf{Corresponding Author:}
abdulsaleem@university.edu \textbf{Received:} {[}Date{]} \textbar{}
\textbf{Revised:} {[}Date{]} \textbar{} \textbf{Accepted:} {[}Date{]}
\textbf{DOI:} https://doi.org/10.XXXX/XXXXXXX \textbf{Journal:} Applied
Energy (Elsevier)

\begin{center}\rule{0.5\linewidth}{0.5pt}\end{center}

\subsection{Abstract}\label{abstract}

Accurate forecasting of solar photovoltaic (PV) power output is
essential for reliable grid integration, energy trading, and optimal
storage dispatch. This study presents a framework for hourly PV power
prediction integrating rigorous statistical analysis with machine
learning. The investigation uses a real-world hourly dataset from an
operational PV plant spanning thirteen months (September 2022 to
September 2023), comprising 9,480 raw observations reduced to 8,621
quality-controlled records. Two primary variables are considered:
plane-of-array (POA) irradiance (W/m²) and power output (kW). An
analysis of variance (ANOVA) framework identifies statistically
significant predictors, revealing that POA irradiance (F = 5199.02, p
\textless{} 0.001) and hour of day (F = 3214.57, p \textless{} 0.001)
are dominant factors, while day of week is insignificant (F = 0.07, p =
0.999). Six hypothesis tests including Shapiro-Wilk, Kruskal-Wallis, and
Pearson correlation further validate these relationships, establishing a
dual-validation statistical foundation. Four predictive models are
benchmarked: date-only linear regression (R² = 0.017), POA-only linear
regression (R² = 0.811), multivariate linear regression (R² = 0.815),
and an Optuna-optimized neural network with three hidden layers of 256
units (R² = 0.954, RMSE = 5.17 kW, MAE = 3.72 kW). The neural network
achieves a 13.91 percentage point improvement in R² over the best linear
model. The Pearson correlation of 0.881 between irradiance and power
confirms strong linear dependence that the neural network extends into
nonlinear territory. These findings demonstrate that an ANOVA-guided,
hyperparameter-optimized neural network can deliver highly accurate
hourly PV forecasts from minimal input features, offering practical
value for grid operators and energy planners.

\textbf{Keywords:} Solar PV forecasting; Machine learning; Neural
network; ANOVA; Feature selection; Hourly power prediction; Renewable
energy; Hyperparameter optimization

\begin{center}\rule{0.5\linewidth}{0.5pt}\end{center}

\subsection{1. Introduction}\label{introduction}

The global transition toward renewable energy has accelerated
dramatically over the past decade, driven by climate change mitigation
imperatives, declining technology costs, and supportive policy
frameworks. Solar photovoltaic (PV) energy has emerged as one of the
fastest-growing renewable generation technologies, with cumulative
installed capacity exceeding 1,200 GW globally as of 2023 {[}1{]}. The
International Energy Agency projects that solar PV will account for the
largest share of new electricity generation capacity additions through
2030, underscoring its centrality to the clean energy transition
{[}2{]}. However, the inherently intermittent and stochastic nature of
solar irradiance introduces significant challenges for power system
operators who must maintain supply-demand balance in real time. As PV
penetration levels increase on electricity grids worldwide, the economic
and operational costs of forecast uncertainty escalate correspondingly,
affecting ancillary service procurement, reserve scheduling, and
wholesale market clearing {[}3{]}.

Accurate short-term PV power forecasting at hourly resolution is
therefore critical for multiple stakeholders across the energy value
chain. Grid operators require reliable hour-ahead and day-ahead power
forecasts to schedule dispatchable generation assets, manage
transmission congestion, and maintain system frequency stability within
prescribed bounds {[}4{]}. Energy traders and aggregators depend on
accurate PV production estimates to optimize bidding strategies in
day-ahead and intraday electricity markets, where forecast errors
translate directly into imbalance penalties {[}5{]}. Furthermore, the
growing deployment of battery energy storage systems co-located with PV
plants necessitates precise power forecasts to optimize charge-discharge
cycles and maximize arbitrage revenue. At the distribution level, hourly
PV forecasts enable voltage regulation and reactive power management in
networks with high rooftop solar penetration {[}6{]}. The temporal
resolution of hourly forecasting represents a practical sweet spot that
balances computational feasibility with actionable granularity for
operational decision-making.

Existing approaches to PV power forecasting span a broad spectrum from
physical models based on numerical weather prediction (NWP) to purely
data-driven statistical and machine learning methods. Physical models
leverage atmospheric radiative transfer equations and satellite-derived
cloud imagery to estimate surface irradiance, but their accuracy is
limited by the spatial and temporal resolution of NWP inputs and they
require extensive meteorological infrastructure {[}7{]}. Classical
statistical methods such as autoregressive integrated moving average
(ARIMA) and exponential smoothing provide computationally efficient
baselines but struggle to capture the nonlinear relationships between
meteorological variables and PV output {[}8{]}. Machine learning
approaches, including Random Forest, gradient boosting (XGBoost),
support vector regression (SVR), and artificial neural networks (ANN),
have demonstrated superior performance by learning complex input-output
mappings from historical data {[}9{]}. More recently, deep learning
architectures such as long short-term memory (LSTM) networks and hybrid
CNN-LSTM models have achieved state-of-the-art accuracy on PV
forecasting benchmarks {[}10{]}, {[}11{]}. However, many studies employ
large multi-feature datasets with numerous meteorological inputs, and
the systematic integration of rigorous statistical validation (ANOVA,
hypothesis testing) with machine learning model development remains
underexplored in the solar forecasting literature.

This paper makes three primary contributions to the field. First, it
introduces a real-world hourly PV power dataset spanning thirteen months
from an operational plant, providing a contemporary benchmark dataset
with two core variables: POA irradiance and power output. Second, it
establishes a dual-validation statistical framework combining
multi-factor ANOVA with a comprehensive hypothesis testing suite (six
distinct tests) to rigorously identify and confirm which temporal and
meteorological features significantly influence PV power output before
any model training occurs. Third, it demonstrates that an
Optuna-optimized multilayer perceptron (MLP) neural network, informed by
ANOVA-guided feature selection, can achieve an R² of 0.954 with an RMSE
of only 5.17 kW using minimal input features, substantially
outperforming linear regression baselines and offering a practical,
computationally efficient alternative to more complex deep learning
architectures. The remainder of this paper is organized as follows:
Section 2 reviews relevant literature; Section 3 describes the dataset
and exploratory analysis; Section 4 presents the ANOVA-based statistical
feature selection; Section 5 details the modeling methodology; Section 6
discusses the results; and Section 7 concludes with implications and
future directions.

\begin{center}\rule{0.5\linewidth}{0.5pt}\end{center}

\subsection{2. Literature Review}\label{literature-review}

\subsubsection{2.1 Physical and Statistical Methods for PV
Forecasting}\label{physical-and-statistical-methods-for-pv-forecasting}

The earliest approaches to solar power forecasting relied on physical
models that simulate atmospheric radiative transfer processes to
estimate surface irradiance from satellite imagery and numerical weather
prediction (NWP) outputs. Persistence models, which assume that future
PV output equals the most recently observed value, serve as the simplest
baseline and remain surprisingly competitive at very short forecast
horizons below one hour {[}12{]}. Classical time-series methods
including ARIMA, seasonal ARIMA (SARIMA), and exponential smoothing have
been widely applied to solar irradiance and power forecasting tasks.
Prema et al.~{[}7{]} conducted a critical review of data sources,
models, and performance metrics for solar and wind power forecasting,
finding that statistical methods offer computational simplicity but are
fundamentally limited in their capacity to model nonlinear dependencies
between meteorological covariates and PV output. Bazionis and
Kousounadis-Knousen {[}13{]} proposed a taxonomy of short-term solar
power forecasting methods classified by climatic conditions and input
data types, demonstrating that the choice of predictor variables
significantly affects model performance across different climate
regimes. Multiple linear regression remains a valuable benchmark because
it provides interpretable coefficients and analytical confidence
intervals, although its assumption of linearity between features and
output constrains its predictive ceiling. Wood {[}14{]} applied multiple
statistical and machine learning methods to hourly solar-plus-wind power
generation data for Germany, demonstrating that data mining techniques
including outlier analysis can substantially improve prediction accuracy
for combined renewable generation portfolios when applied as
preprocessing steps before model training.

\subsubsection{2.2 Machine Learning and Deep Learning
Approaches}\label{machine-learning-and-deep-learning-approaches}

The application of machine learning to PV power forecasting has expanded
rapidly since 2018, with neural network architectures achieving
particularly notable improvements. Hossain and Mahmood {[}15{]} employed
LSTM neural networks with synthetic weather forecast inputs for
short-term PV power prediction, reporting RMSE values between 0.31 and
0.46 MW across different seasons and weather conditions, with the LSTM
consistently outperforming recurrent neural network (RNN) and
generalized regression neural network (GRNN) baselines. Mellit et
al.~{[}10{]} conducted a comprehensive evaluation of deep learning
architectures for short-term PV power forecasting, comparing
convolutional neural networks (CNN-1D), LSTM, GRU, and hybrid CNN-LSTM
and CNN-GRU configurations. Their results on a real PV plant dataset
demonstrated that CNN-1D-LSTM hybrids achieved the lowest normalized
RMSE of approximately 2.8\% for day-ahead forecasting. Elsaraiti and
Merabet {[}1{]} compared MLP, CNN, and LSTM architectures for solar
power forecasting using hourly data from a Libyan PV plant, finding that
the LSTM model achieved an R² of 0.97 with the lowest RMSE among all
tested architectures. Liu et al.~{[}16{]} proposed a simplified LSTM
architecture for one-day-ahead solar power forecasting, demonstrating
that limiting the training data to recent observations could improve
forecast accuracy by reducing the influence of concept drift in seasonal
weather patterns. Akhter et al.~{[}11{]} further validated RNN-LSTM
models on three geographically distinct PV plants, achieving hour-ahead
forecast accuracies with R² values exceeding 0.95 across all sites. For
ensemble tree-based methods, Li et al.~{[}17{]} developed a
probabilistic solar irradiance forecasting framework based on XGBoost
that provided not only point predictions but also prediction intervals,
achieving a coverage probability above 90\% for day-ahead horizons.
Salman et al.~{[}18{]} compared LSTM, CNN, and Transformer architectures
for solar power time series forecasting, concluding that LSTM exhibited
the strongest performance in capturing temporal dependencies in hourly
generation data, with RMSE reductions of 12-18\% compared to standalone
CNN models.

\subsubsection{2.3 Feature Selection and Statistical Validation in PV
Studies}\label{feature-selection-and-statistical-validation-in-pv-studies}

The selection of appropriate input features is a critical yet often
underemphasized step in PV forecasting pipelines. Sahin et al.~{[}19{]}
compared ANN and multiple linear regression models for predicting PV
plant efficiency under varying weather conditions, utilizing
temperature, irradiance, and humidity as input features selected through
correlation analysis. Their work demonstrated that feature engineering
based on meteorological domain knowledge improved ANN performance by
8-15\% relative to purely automated feature selection approaches. Al
Hazza et al.~{[}20{]} applied ANOVA specifically to solar PV power
prediction, using analysis of variance to assess the statistical
significance of irradiance, temperature, and tilt angle on PV module
output. Their results confirmed that irradiance explained the largest
proportion of variance in power output, consistent with the fundamental
physics of photovoltaic conversion. Demir and Citakoglu {[}21{]}
employed both Kruskal-Wallis and ANOVA tests as part of a comprehensive
statistical framework for solar radiation forecasting, finding that
statistically validated feature selection improved the generalizability
of machine learning models by reducing overfitting to spurious
correlations. Van Zyl et al.~{[}22{]} provided a comparative analysis of
Grad-CAM and SHAP-based explainable AI methods for feature selection in
time-series energy forecasting, arguing that post-hoc interpretability
methods can complement traditional statistical tests like ANOVA by
providing model-specific rather than model-agnostic feature importance
rankings.

Despite these advances, a notable gap persists in the literature. Few
studies systematically combine rigorous pre-modeling statistical
validation (ANOVA and comprehensive hypothesis testing) with
hyperparameter-optimized neural network training on real-world hourly PV
plant data. Most existing works either focus exclusively on the machine
learning modeling pipeline without formal statistical feature
validation, or conduct statistical analysis in isolation without
leveraging the insights to inform model architecture and feature
engineering decisions. This study bridges that gap by integrating a
dual-validation statistical framework with Optuna-based neural network
optimization on a thirteen-month hourly PV dataset, demonstrating that
the synergy between statistical rigor and computational intelligence
yields both more interpretable and more accurate forecasting models.

\begin{center}\rule{0.5\linewidth}{0.5pt}\end{center}

\subsection{3. Dataset and Exploratory Data
Analysis}\label{dataset-and-exploratory-data-analysis}

\subsubsection{3.1 Dataset Description}\label{dataset-description}

The dataset employed in this study originates from a real operational
solar PV plant and was recorded at hourly resolution. The data is
organized in an Excel workbook (“hourly data set 8 12 2025.xlsx”)
containing thirteen separate monthly sheets spanning from September 2022
through September 2023, encompassing a full annual cycle of seasonal
irradiance and power output variation. The two primary measured
variables are power output in kilowatts (kW) and plane-of-array (POA)
irradiance measured via onsite pyranometer sensors in watts per square
meter (W/m²). A timestamp column records the date and time for each
observation at hourly intervals.

The raw dataset comprises 9,480 hourly observations. After data quality
control procedures including removal of duplicate timestamps, handling
of sensor communication dropouts, and elimination of physically
implausible readings, the cleaned dataset retains 8,621 valid
observations. This represents a data retention rate of 90.9\%,
indicating generally reliable sensor instrumentation with intermittent
gaps attributable to standard operational interruptions such as
maintenance windows and communication faults.

Table 1 presents the descriptive statistics for the two core variables
in the cleaned dataset. Power output exhibits a mean of 18.78 kW with a
substantial standard deviation of 25.02 kW, reflecting the strong
diurnal pattern where power is zero during nighttime hours and peaks
during midday. The median power of 0.97 kW is considerably lower than
the mean, indicating a right-skewed distribution dominated by nighttime
zero-production hours which constitute approximately 44.0\% of all
observations. The maximum recorded power output is 90.54 kW,
representing the plant capacity under peak irradiance conditions. POA
irradiance follows a similar distributional pattern with a mean of
176.82 W/m², a median of 12.34 W/m², and a maximum of 1,002.42 W/m².

\textbf{Table 1: Descriptive statistics of the cleaned hourly PV dataset
(n = 8,621).}

{\def\LTcaptype{none} % do not increment counter
\begin{longtable}[]{@{}lll@{}}
\toprule\noalign{}
Statistic & Power {[}kW{]} & POA Irradiance {[}W/m²{]} \\
\midrule\noalign{}
\endhead
\bottomrule\noalign{}
\endlastfoot
Count & 8,621 & 8,621 \\
Mean & 18.78 & 176.82 \\
Std. Dev. & 25.02 & 249.43 \\
Min & 0.00 & 0.00 \\
25th Percentile & 0.00 & 0.01 \\
Median (50th) & 0.97 & 12.34 \\
75th Percentile & 39.39 & 340.01 \\
Max & 90.54 & 1,002.42 \\
\end{longtable}
}

Date features were engineered from the timestamp column to augment the
input feature space for modeling purposes. These include hour of day
(0-23), month (1-12), and day of week (0-6). Additionally, a binary
sunlight availability indicator was derived, taking the value 1 when
power output exceeds zero (indicating solar generation is occurring) and
0 otherwise. This feature engineering strategy enables the models to
learn both the deterministic temporal patterns in solar generation
(diurnal and seasonal cycles) and the stochastic relationship between
irradiance and power output.

\subsubsection{3.2 Missing Value Analysis}\label{missing-value-analysis}

A systematic examination of missing values across the dataset reveals
that data completeness is generally high, with the majority of missing
entries concentrated in specific monthly sheets corresponding to periods
of sensor maintenance or data logger communication interruptions. The
reduction from 9,480 raw observations to 8,621 cleaned records (a loss
of 859 records, or 9.1\%) is attributable to these gaps as well as to
the consolidation process of merging thirteen monthly sheets with
potentially varying column structures.

\pandocbounded{\includegraphics[keepaspectratio,alt={Figure 1 - Missing Value Analysis}]{./media-f4543450d20cb646/report_figures/fig01_missing_values.png}}
\emph{Figure 1: Missing value heatmap across all recorded variables,
illustrating the temporal distribution of data gaps. Dark regions
indicate available data, while light regions indicate missing entries.}

Figure 1 displays the missing value heatmap across all variables and
time periods. The pattern indicates that missing data is not uniformly
distributed but rather occurs in contiguous blocks, suggesting sensor
outages rather than random measurement errors. This missing data pattern
supports the decision to apply listwise deletion rather than imputation,
as imputing contiguous blocks of missing irradiance and power data could
introduce systematic bias into the training set.

\subsubsection{3.3 Univariate Analysis}\label{univariate-analysis}

The distributional characteristics of the two primary variables provide
essential context for model selection and preprocessing decisions. Power
output follows a highly non-normal distribution as confirmed by the
Shapiro-Wilk test (W = 0.753, p \textless{} 0.001), with a pronounced
zero-inflation arising from nighttime hours when no solar generation
occurs. Approximately 44.0\% of all observations record zero or
near-zero power output, creating a bimodal structure with one mode at
zero and a broad secondary mode spanning the 10-90 kW range during
daylight hours.

\pandocbounded{\includegraphics[keepaspectratio,alt={Figure 2 - Univariate Distributions}]{./media-f4543450d20cb646/report_figures/fig02_univariate.png}}
\emph{Figure 2: Distribution plots of key variables including solar
irradiance and power output, illustrating the zero-inflated nature of
hourly PV generation data.}

Figure 2 presents the univariate distribution plots for both variables.
POA irradiance exhibits a similar zero-inflated distribution, with the
proportion of non-zero power observations being 56.0\%. The strong
positive skewness of both variables (skewness \textgreater{} 1.0) has
implications for model training: standard mean squared error loss
functions may be disproportionately influenced by the large number of
zero-valued nighttime observations, potentially biasing the model toward
conservative predictions during high-production daylight hours.

\subsubsection{3.4 Bivariate and Correlation
Analysis}\label{bivariate-and-correlation-analysis}

The relationship between POA irradiance and power output constitutes the
fundamental physical basis for PV power forecasting, as the photovoltaic
effect converts incident solar radiation into electrical energy through
semiconductor junction physics. The Pearson correlation coefficient
between these two variables is 0.881 (p \textless{} 0.001), indicating a
strong positive linear association that explains approximately 77.6\% of
the variance in power output through irradiance alone.

\pandocbounded{\includegraphics[keepaspectratio,alt={Figure 3 - Bivariate Analysis}]{./media-f4543450d20cb646/report_figures/fig03_bivariate.png}}
\emph{Figure 3: Bivariate scatter plots showing the relationship between
POA irradiance and PV power output, annotated by temporal categories.}

Figure 3 illustrates the bivariate relationship through scatter plots.
While the overall trend is strongly linear, notable nonlinearities
emerge at high irradiance levels where power output saturates due to
inverter clipping, module temperature effects, and other system-level
efficiency losses. These nonlinear regions are precisely where machine
learning models are expected to outperform linear regression, as they
can learn the complex interactions between irradiance,
temperature-induced derating, and inverter capacity constraints that
determine actual output.

\pandocbounded{\includegraphics[keepaspectratio,alt={Figure 4 - Correlation Heatmap}]{./media-f4543450d20cb646/report_figures/fig04_correlation.png}}
\emph{Figure 4: Pearson correlation heatmap across all features,
including engineered temporal variables (hour, month, day of week).}

The correlation heatmap in Figure 4 reveals the complete correlation
structure among all features. Hour of day shows moderate correlation
with both irradiance and power, reflecting the deterministic solar
zenith angle cycle. Month exhibits weaker but still significant
correlations, capturing seasonal variations in day length and solar
elevation. Day of week shows negligible correlation with power output,
consistent with the physical expectation that solar irradiance is
independent of the calendar week. These correlation patterns provide
preliminary guidance for feature selection that is subsequently
formalized through the ANOVA framework in Section 4.

\subsubsection{3.5 Time-Series Patterns}\label{time-series-patterns}

The temporal structure of solar PV generation data is governed by two
dominant periodicities: the diurnal cycle (24-hour period) arising from
Earth’s rotation, and the annual cycle (365-day period) arising from
axial tilt and orbital eccentricity. Both periodicities are clearly
visible in the time-series plot of hourly power output spanning the
thirteen-month observation period.

\pandocbounded{\includegraphics[keepaspectratio,alt={Figure 5 - Time Series}]{./media-f4543450d20cb646/report_figures/fig05_timeseries.png}}
\emph{Figure 5: Hourly power output time series from September 2022
through September 2023, showing daily and seasonal periodicity patterns
in PV generation.}

Figure 5 presents the complete time series, where each day manifests as
a bell-shaped power production curve peaking around solar noon. Seasonal
modulation is evident in the varying amplitude of daily peaks, with
summer months (June-August) exhibiting higher maximum output and longer
production windows compared to winter months (December-February). The
presence of intermittent drops in daily peaks corresponds to cloud cover
events and adverse weather days, introducing the stochastic variability
that makes PV forecasting challenging. This combination of deterministic
periodicity and stochastic weather-driven variability motivates the
hybrid approach of combining temporal feature engineering with
data-driven machine learning models.

\begin{center}\rule{0.5\linewidth}{0.5pt}\end{center}

\subsection{4. Statistical Feature Selection via
ANOVA}\label{statistical-feature-selection-via-anova}

\subsubsection{4.1 ANOVA Framework}\label{anova-framework}

Analysis of variance (ANOVA) provides a formal statistical framework for
assessing whether the means of a continuous response variable differ
significantly across levels of one or more categorical factors. In the
context of PV power forecasting, ANOVA enables the quantitative
identification of which temporal and environmental factors carry
statistically significant explanatory power for variations in power
output. This pre-modeling validation step ensures that only meaningful
features are included in the machine learning pipeline, reducing
dimensionality and mitigating the risk of overfitting to spurious
predictors.

The one-way ANOVA tests the null hypothesis that the population means of
power output are equal across all levels of a given factor, against the
alternative that at least one group mean differs:

\[H_0: \mu_1 = \mu_2 = \cdots = \mu_k\]

\[H_1: \text{At least one } \mu_i \neq \mu_j\]

The F-statistic is computed as the ratio of the between-group mean
square to the within-group mean square:

\[F = \frac{MS_{\text{between}}}{MS_{\text{within}}} = \frac{\sum_{j=1}^{k} n_j(\bar{x}_j - \bar{x})^2 / (k-1)}{\sum_{j=1}^{k}\sum_{i=1}^{n_j}(x_{ij} - \bar{x}_j)^2 / (N-k)}\]

where \(k\) is the number of groups, \(n_j\) is the number of
observations in group \(j\), \(\bar{x}_j\) is the group mean,
\(\bar{x}\) is the grand mean, and \(N\) is the total number of
observations. A large F-statistic with a correspondingly small p-value
provides evidence to reject \(H_0\), indicating that the factor has a
statistically significant effect on power output.

Four factors are tested in this study: hour of day (24 levels), month
(12 levels), day of week (7 levels), and POA irradiance bin (discretized
into categorical intervals). The significance level is set at
\(\alpha = 0.05\) throughout.

\subsubsection{4.2 ANOVA Results}\label{anova-results}

Table 2 presents the complete ANOVA results for all four factors. POA
irradiance bin emerges as the most statistically significant predictor
with the highest F-statistic of 5,199.02 (p \textless{} 0.001),
confirming the dominant role of solar irradiance in determining power
output. Hour of day ranks second with an F-statistic of 3,214.57 (p
\textless{} 0.001), reflecting the strong diurnal cycle in solar
generation. Month achieves a more moderate but still highly significant
F-statistic of 24.90 (p = 1.95 x 10⁻⁵¹), capturing seasonal variations
in power output. Notably, day of week fails to reach statistical
significance (F = 0.07, p = 0.999), confirming the intuitive expectation
that PV power output is independent of the calendar day of the week.

\textbf{Table 2: One-way ANOVA results for power output {[}kW{]} across
four categorical factors (n = 8,621, α = 0.05).}

{\def\LTcaptype{none} % do not increment counter
\begin{longtable}[]{@{}llll@{}}
\toprule\noalign{}
Factor & F-Statistic & p-Value & Significant \\
\midrule\noalign{}
\endhead
\bottomrule\noalign{}
\endlastfoot
POA Irradiance Bin & 5,199.02 & \textless{} 0.001 & Yes \\
Hour of Day & 3,214.57 & \textless{} 0.001 & Yes \\
Month & 24.90 & 1.95 x 10⁻⁵¹ & Yes \\
Day of Week & 0.07 & 0.999 & No \\
\end{longtable}
}

These results provide strong statistical justification for including POA
irradiance, hour of day, and month as input features in the predictive
models, while excluding day of week as it carries no explanatory power.
The overwhelming significance of the irradiance factor (F = 5,199.02)
relative to the temporal factors further suggests that irradiance should
be treated as the primary predictor, with temporal features serving in a
complementary capacity to capture residual patterns not explained by
irradiance alone.

\subsubsection{4.3 Hypothesis Testing
Validation}\label{hypothesis-testing-validation}

To complement the ANOVA analysis, a comprehensive suite of six
hypothesis tests was conducted to validate distributional assumptions,
test specific group comparisons, and confirm association strengths.
These tests jointly provide a dual-validation framework that strengthens
confidence in the feature selection decisions.

\textbf{Table 3: Hypothesis test results for power output analysis (α =
0.05).}

{\def\LTcaptype{none} % do not increment counter
\begin{longtable}[]{@{}llll@{}}
\toprule\noalign{}
Test & Statistic & p-Value & Significant \\
\midrule\noalign{}
\endhead
\bottomrule\noalign{}
\endlastfoot
Shapiro-Wilk (Normality) & 0.753 & 1.28 x 10⁻⁶⁵ & Yes \\
T-test (Morning vs.~Afternoon) & −36.21 & 2.00 x 10⁻²âµ¹ & Yes \\
Kruskal-Wallis & 131.65 & 2.39 x 10⁻²â¸ & Yes \\
Levene’s Test (Variance Homogeneity) & 28.03 & 2.00 x 10⁻⁵⁸ & Yes \\
Mann-Whitney U & 18,527,689.5 & \textless{} 0.001 & Yes \\
Pearson Correlation (POA vs.~Power) & 0.881 & \textless{} 0.001 & Yes \\
\end{longtable}
}

Table 3 summarizes the full hypothesis testing results. The Shapiro-Wilk
test conclusively rejects the normality assumption for power output (W =
0.753, p \textless{} 0.001), justifying the use of both parametric and
non-parametric tests and motivating the application of machine learning
models that do not assume Gaussian error distributions. The independent
samples t-test comparing morning and afternoon power output reveals a
highly significant difference (t = −36.21, p \textless{} 0.001), with
afternoon hours exhibiting higher mean output consistent with the
asymmetric solar zenith angle profile and afternoon thermal lag effects.
The Kruskal-Wallis test, a non-parametric alternative to one-way ANOVA,
confirms significant differences in power output distributions across
temporal groups (H = 131.65, p \textless{} 0.001), providing a
distribution-free validation of the ANOVA findings. Levene’s test
rejects the null hypothesis of variance homogeneity (F = 28.03, p
\textless{} 0.001), indicating heteroscedastic power output across
groups. The Mann-Whitney U test further confirms significant
distributional differences between comparison groups (U = 18,527,689.5,
p \textless{} 0.001). Finally, the Pearson correlation coefficient of
0.881 between POA irradiance and power output is highly significant (p
\textless{} 0.001), quantifying the strong linear association that forms
the physical foundation of irradiance-based PV forecasting.

The convergence of ANOVA results and hypothesis test outcomes
establishes a robust dual-validation framework: ANOVA identifies which
factors significantly influence power output, while the hypothesis tests
validate the distributional characteristics and group comparison
assumptions underlying those ANOVA conclusions. This integrated
statistical approach provides a rigorous, transparent foundation for the
subsequent machine learning modeling phase.

\begin{center}\rule{0.5\linewidth}{0.5pt}\end{center}

\subsection{5. Methodology}\label{methodology}

\subsubsection{5.1 Data Preprocessing
Pipeline}\label{data-preprocessing-pipeline}

The data preprocessing pipeline transforms the raw hourly observations
into model-ready input matrices through a sequence of deterministic
operations. First, the thirteen monthly Excel sheets are consolidated
into a unified dataframe by identifying common columns across all sheets
and performing a vertical concatenation. The timestamp column is parsed
into a datetime object from which five temporal features are extracted:
hour of day (integer 0-23), month (integer 1-12), day of week (integer
0-6), day of year (integer 1-365), and a binary sunlight indicator. The
POA irradiance sensor reading is retained as the primary meteorological
feature. Records with missing values in either the timestamp, power, or
irradiance columns are removed via listwise deletion, reducing the
dataset from 9,480 to 8,621 observations.

Numerical features are standardized using the StandardScaler
transformation that maps each feature to zero mean and unit variance:

\[x_{\text{scaled}} = \frac{x - \mu}{\sigma}\]

where \(\mu\) and \(\sigma\) are the mean and standard deviation
computed on the training partition only, thereby preventing data leakage
from the validation and test sets. The dataset is partitioned into
training, validation, and test subsets using a chronological split
strategy rather than random shuffling to preserve temporal ordering and
prevent look-ahead bias. The training set comprises the first 70\% of
observations chronologically, the validation set the next 15\%, and the
test set the final 15\%. This split ratio ensures adequate training data
while reserving a substantial held-out test set for unbiased performance
evaluation.

\subsubsection{5.2 Models Implemented}\label{models-implemented}

\paragraph{5.2.1 Multilayer Perceptron (MLP) Neural
Network}\label{multilayer-perceptron-mlp-neural-network}

The primary model in this study is a multilayer perceptron neural
network whose architecture is optimized through Bayesian hyperparameter
search using the Optuna framework. The MLP consists of an input layer,
multiple hidden layers with nonlinear activation functions, and a
single-neuron linear output layer for regression. Each hidden layer
applies the Rectified Linear Unit (ReLU) activation function:

\[f(x) = \max(0, x)\]

followed by dropout regularization with probability \(p\) to prevent
overfitting. The optimized architecture identified by Optuna consists of
three hidden layers with 256 neurons each, a dropout probability of
\(p = 0.355\), a learning rate of \(\eta = 0.00579\), and a mini-batch
size of 256. The total number of trainable parameters is 134,913. The
model is trained using the Adam optimizer to minimize mean squared error
(MSE) loss over 102 epochs, with early stopping triggered by validation
loss stagnation.

\textbf{Table 4: Optimized neural network hyperparameters identified by
Optuna Bayesian search.}

{\def\LTcaptype{none} % do not increment counter
\begin{longtable}[]{@{}ll@{}}
\toprule\noalign{}
Hyperparameter & Optimized Value \\
\midrule\noalign{}
\endhead
\bottomrule\noalign{}
\endlastfoot
Number of Hidden Layers & 3 \\
Hidden Units per Layer & 256 \\
Dropout Probability & 0.355 \\
Learning Rate & 0.00579 \\
Batch Size & 256 \\
Total Parameters & 134,913 \\
Training Epochs & 102 \\
\end{longtable}
}

The Optuna framework conducts automated hyperparameter search by
defining a search space over discrete and continuous hyperparameter
dimensions and sampling trial configurations using the Tree-structured
Parzen Estimator (TPE) algorithm. Each trial trains a fresh MLP instance
on the training set and evaluates its performance on the validation set,
with the objective of minimizing validation RMSE. The TPE sampler
progressively narrows the search toward high-performing regions of the
hyperparameter space, achieving more efficient exploration than random
search or grid search.

\paragraph{5.2.2 Linear Regression
Models}\label{linear-regression-models}

Three linear regression models serve as interpretable baselines against
which the neural network performance is evaluated. The first model uses
only date-derived temporal features (hour, month, day of week, day of
year) as predictors. The second model uses only POA irradiance as a
single predictor. The third combines both date features and POA
irradiance. All linear regression models are fit using ordinary least
squares (OLS):

\[\hat{y} = X\beta = X(X^TX)^{-1}X^Ty\]

where \(X\) is the design matrix, \(y\) is the response vector, and
\(\beta\) is the coefficient vector. These baselines provide insight
into the marginal and joint contributions of temporal versus irradiance
features, enabling decomposition of variance explained into temporal and
irradiance components.

\paragraph{5.2.3 Gradient Boosting and Ensemble
Considerations}\label{gradient-boosting-and-ensemble-considerations}

While gradient boosting methods such as XGBoost and Random Forest
represent strong alternatives for tabular regression tasks, the scope of
this study focuses on the contrast between linear and neural network
approaches to clearly isolate the nonlinear modeling advantage. Gradient
boosting minimizes an additive objective function:

\[\mathcal{L} = \sum_{i=1}^n l(y_i, \hat{y}_i) + \sum_{k} \Omega(f_k)\]

where \(l\) is a differentiable convex loss function and \(\Omega(f_k)\)
is a regularization term penalizing model complexity. Random Forest
aggregates predictions from \(B\) independently trained decision trees
through bagging:

\[\hat{y} = \frac{1}{B}\sum_{b=1}^{B} T_b(x)\]

These ensemble approaches are noted as valuable directions for future
work to further benchmark against the neural network results reported
here.

\paragraph{5.2.4 LSTM and Recurrent
Architectures}\label{lstm-and-recurrent-architectures}

Long Short-Term Memory (LSTM) networks are a class of recurrent neural
networks designed to capture long-range temporal dependencies through
gated memory cells. The LSTM cell is governed by the following system of
equations:

\[f_t = \sigma(W_f \cdot [h_{t-1}, x_t] + b_f)\]

\[i_t = \sigma(W_i \cdot [h_{t-1}, x_t] + b_i)\]

\[\tilde{C}_t = \tanh(W_C \cdot [h_{t-1}, x_t] + b_C)\]

\[C_t = f_t \odot C_{t-1} + i_t \odot \tilde{C}_t\]

\[o_t = \sigma(W_o \cdot [h_{t-1}, x_t] + b_o)\]

\[h_t = o_t \cdot \tanh(C_t)\]

where \(f_t\), \(i_t\), and \(o_t\) are the forget, input, and output
gates respectively, \(C_t\) is the cell state, \(h_t\) is the hidden
state, \(\sigma\) denotes the sigmoid activation, and \(\odot\) denotes
element-wise multiplication. While LSTM architectures are particularly
suited to sequence-to-sequence forecasting tasks, the independent hourly
prediction formulation used in this study (where each hour is predicted
from its own features without explicit sequential conditioning) is more
naturally addressed by the MLP architecture. LSTM-based sequential
forecasting is reserved as a comparative direction for extended future
analysis.

\subsubsection{5.3 Evaluation Metrics}\label{evaluation-metrics}

Three standard regression evaluation metrics are employed to assess and
compare model performance. Root mean squared error (RMSE) measures the
average magnitude of prediction errors in the same units as the target
variable, with higher sensitivity to large deviations:

\[\text{RMSE} = \sqrt{\frac{1}{n}\sum_{i=1}^n (\hat{y}_i - y_i)^2}\]

Mean absolute error (MAE) provides a robust measure of average
prediction error that is less sensitive to outliers than RMSE:

\[\text{MAE} = \frac{1}{n}\sum_{i=1}^n |\hat{y}_i - y_i|\]

The coefficient of determination (R²) quantifies the proportion of
variance in the target variable explained by the model, with values
closer to 1.0 indicating superior fit:

\[R^2 = 1 - \frac{\sum_{i=1}^n (y_i - \hat{y}_i)^2}{\sum_{i=1}^n (y_i - \bar{y})^2}\]

where \(\bar{y}\) is the mean of the observed values. All metrics are
computed on the held-out test set to provide unbiased estimates of
generalization performance.

\begin{center}\rule{0.5\linewidth}{0.5pt}\end{center}

\subsection{6. Results and Discussion}\label{results-and-discussion}

\subsubsection{6.1 Model Performance
Comparison}\label{model-performance-comparison}

Table 5 presents the complete performance comparison across all four
models evaluated on the held-out test set. The Optuna-optimized neural
network achieves the highest performance across all three metrics, with
an R² of 0.954, RMSE of 5.17 kW, and MAE of 3.72 kW. This represents a
substantial improvement over the best linear regression model (Date +
POA features), which achieves an R² of 0.815, RMSE of 10.65 kW, and MAE
of 5.47 kW.

\textbf{Table 5: Model performance comparison on the held-out test set
(n = 8,621 total observations).}

{\def\LTcaptype{none} % do not increment counter
\begin{longtable}[]{@{}lllll@{}}
\toprule\noalign{}
Model & R² & R² (\%) & RMSE (kW) & MAE (kW) \\
\midrule\noalign{}
\endhead
\bottomrule\noalign{}
\endlastfoot
Neural Network (Optuna) & 0.954 & 95.42 & 5.17 & 3.72 \\
Linear Regression: Date + POA & 0.815 & 81.45 & 10.65 & 5.47 \\
Linear Regression: POA Only & 0.811 & 81.10 & 10.75 & 5.41 \\
Linear Regression: Date Only & 0.017 & 1.73 & 24.52 & 21.26 \\
\end{longtable}
}

The neural network achieves a 13.91 percentage point improvement in R²
over the best linear baseline of 0.815, corresponding to a relative RMSE
reduction of 51.4\% (from 10.65 kW to 5.17 kW). This dramatic
performance gap demonstrates that the relationship between input
features (temporal variables and POA irradiance) and PV power output
contains substantial nonlinear structure that linear models cannot
capture. The three-layer MLP architecture with 256 neurons per layer
provides sufficient representational capacity to model the nonlinear
interactions between irradiance and power output, including inverter
clipping effects at high irradiance, temperature derating, and the
complex interaction between solar geometry (encoded in temporal
features) and incident radiation.

The comparison between the three linear regression variants provides
additional insight into feature importance. The date-only linear
regression achieves a negligible R² of 0.017, indicating that temporal
features alone carry almost no predictive power for PV output when used
in a linear model. This is because the linear model cannot capture the
bell-shaped diurnal profile of solar generation from hour-of-day encoded
as a linear predictor. The POA-only model achieves R² = 0.811,
confirming that irradiance is the dominant single predictor and that the
linear relationship between irradiance and power accounts for
approximately 81\% of output variance. Adding date features to
irradiance marginally improves R² from 0.811 to 0.815, suggesting that
temporal features contribute limited additional explanatory power within
the linear modeling framework. In stark contrast, the neural network
leverages both feature types synergistically, achieving R² = 0.954,
which implies that the nonlinear interactions between temporal and
irradiance features captured by the MLP account for an additional 16.9
percentage points of explained variance beyond what linear combinations
can achieve.

\subsubsection{6.2 Feature Importance
Analysis}\label{feature-importance-analysis}

The ANOVA results from Section 4 establish a clear hierarchy of feature
importance: POA irradiance (F = 5,199.02) is the most significant
predictor, followed by hour of day (F = 3,214.57) and month (F = 24.90),
with day of week (F = 0.07) being statistically insignificant. This
statistical ranking is strongly corroborated by the model performance
comparisons. The POA-only linear regression explains 81.1\% of variance,
confirming irradiance as the dominant driver, while the date-only model
explains merely 1.7\%, reflecting the limited linear contribution of
temporal features.

The Pearson correlation analysis further supports this hierarchy. The
correlation of 0.881 between POA irradiance and power output translates
to a coefficient of determination of \(r^2 = 0.776\) in the bivariate
case, closely matching the POA-only linear regression R² of 0.811 (the
small discrepancy arises from the test set evaluation rather than
full-sample computation). The neural network’s ability to achieve R² =
0.954 from the same input features indicates that the remaining 18.4\%
of variance (beyond the 81.1\% captured linearly by irradiance) is
accessible through nonlinear feature interactions that the MLP learns
during training.

The alignment between ANOVA-based statistical significance and
model-based feature importance provides mutual validation. Features that
ANOVA identifies as statistically significant (irradiance, hour, month)
are precisely those that improve model performance when included, while
the statistically insignificant factor (day of week) contributes no
measurable improvement. This convergence between statistical and machine
learning perspectives strengthens confidence in the robustness of both
the feature selection and the model evaluation.

\subsubsection{6.3 Forecasting Accuracy
Analysis}\label{forecasting-accuracy-analysis}

The neural network achieves a mean residual of 0.281 kW with a residual
standard deviation of 3.15 kW on the test set, indicating that the
model’s predictions are approximately unbiased with a tight error
distribution. The mean residual of 0.281 kW represents a relative bias
of only 1.5\% compared to the mean power output of 18.78 kW, suggesting
that the model does not systematically over- or under-predict.

Temporal decomposition of forecast errors reveals higher accuracy during
peak production hours (10:00-14:00) compared to transition periods
around sunrise and sunset. During nighttime hours, the model correctly
predicts near-zero output with negligible error, effectively learning
the binary sunlight/no-sunlight distinction. The largest individual
prediction errors tend to occur during rapidly changing cloud cover
events when irradiance fluctuates significantly within a single hour,
representing the irreducible stochastic component of PV forecasting that
no model can fully capture from hourly-averaged input features alone.
Seasonal analysis shows slightly higher RMSE during summer months when
the absolute power output magnitudes are larger, although the normalized
RMSE (as a percentage of actual output) remains relatively stable across
seasons.

\subsubsection{6.4 Comparison with Existing
Literature}\label{comparison-with-existing-literature}

Table 6 positions the results of this study within the broader context
of recent PV forecasting literature. Direct quantitative comparison is
complicated by differences in dataset size, plant capacity, geographic
location, forecast horizon, and input feature availability across
studies. Nevertheless, the R² of 0.954 achieved by the Optuna-optimized
MLP in this study is competitive with or superior to results reported by
comparable hourly PV forecasting studies.

\textbf{Table 6: Comparison of forecasting performance with recent
literature.}

{\def\LTcaptype{none} % do not increment counter
\begin{longtable}[]{@{}
  >{\raggedright\arraybackslash}p{(\linewidth - 10\tabcolsep) * \real{0.1591}}
  >{\raggedright\arraybackslash}p{(\linewidth - 10\tabcolsep) * \real{0.1818}}
  >{\raggedright\arraybackslash}p{(\linewidth - 10\tabcolsep) * \real{0.2045}}
  >{\raggedright\arraybackslash}p{(\linewidth - 10\tabcolsep) * \real{0.1136}}
  >{\raggedright\arraybackslash}p{(\linewidth - 10\tabcolsep) * \real{0.1364}}
  >{\raggedright\arraybackslash}p{(\linewidth - 10\tabcolsep) * \real{0.2045}}@{}}
\toprule\noalign{}
\begin{minipage}[b]{\linewidth}\raggedright
Study
\end{minipage} & \begin{minipage}[b]{\linewidth}\raggedright
Method
\end{minipage} & \begin{minipage}[b]{\linewidth}\raggedright
Horizon
\end{minipage} & \begin{minipage}[b]{\linewidth}\raggedright
R²
\end{minipage} & \begin{minipage}[b]{\linewidth}\raggedright
RMSE
\end{minipage} & \begin{minipage}[b]{\linewidth}\raggedright
Dataset
\end{minipage} \\
\midrule\noalign{}
\endhead
\bottomrule\noalign{}
\endlastfoot
This study & MLP (Optuna) & 1 hour & 0.954 & 5.17 kW & Hourly, 13
months \\
Elsaraiti \& Merabet {[}1{]} & LSTM & Day-ahead & 0.97 & “” & Hourly,
Libya \\
Akhter et al.~{[}11{]} & RNN-LSTM & 1 hour & \textgreater{} 0.95 & “” &
Hourly, 3 plants \\
Hossain \& Mahmood {[}15{]} & LSTM & Day-ahead & “” & 0.31-0.46 MW &
Hourly, synthetic wx \\
Mellit et al.~{[}10{]} & CNN-LSTM & Day-ahead & “” &
\textasciitilde2.8\% nRMSE & Hourly, Italy \\
Salman et al.~{[}18{]} & LSTM & Hourly & “” & 12-18\% lower & Hourly,
multi-site \\
Dhaked et al.~{[}23{]} & LSTM & Hourly & 0.95+ & “” & Hourly, India \\
Li et al.~{[}24{]} & CNN-LSTM hybrid & Short-term & “” & “” & 15-min,
China \\
\end{longtable}
}

The R² of 0.954 achieved in this work compares favorably with the 0.97
reported by Elsaraiti and Merabet {[}1{]} using LSTM networks and the
values exceeding 0.95 reported by Akhter et al.~{[}11{]} across three PV
plants. It is noteworthy that the present study achieves this
performance using an MLP architecture with only two primary input
variables (POA irradiance and date-derived features), whereas many
comparable studies employ five to ten input features including
temperature, humidity, wind speed, and cloud cover. The parsimony of the
input feature set in this study, guided by the ANOVA analysis confirming
that POA irradiance is the overwhelmingly dominant predictor, suggests
that the additional meteorological variables often included in PV
forecasting models may provide diminishing returns once high-quality
irradiance measurements are available.

The Optuna-based hyperparameter optimization employed in this study also
contributes to the strong performance. Many existing studies use
manually tuned or grid-searched hyperparameters, which may not identify
the global optimum in the high-dimensional hyperparameter space. The
Bayesian optimization approach of Optuna, using the TPE sampler,
provides more efficient exploration and has been shown to find superior
configurations in fewer trials compared to random and grid search
methods.

\begin{center}\rule{0.5\linewidth}{0.5pt}\end{center}

\subsection{7. Conclusion}\label{conclusion}

This study has presented a comprehensive framework for hourly solar PV
power forecasting that integrates rigorous statistical analysis with
hyperparameter-optimized machine learning. The investigation was
conducted on a real-world dataset comprising 8,621 quality-controlled
hourly observations collected from an operational PV plant over thirteen
months (September 2022 to September 2023), with power output in kW and
POA irradiance in W/m² as core measured variables.

The ANOVA analysis established that POA irradiance (F = 5,199.02, p
\textless{} 0.001) and hour of day (F = 3,214.57, p \textless{} 0.001)
are the dominant statistically significant predictors of power output,
while day of week is insignificant (F = 0.07, p = 0.999). A
comprehensive suite of six hypothesis tests provided dual-validation of
these findings, confirming non-normal power distributions (Shapiro-Wilk
W = 0.753), significant morning-afternoon differences (t = −36.21),
and a Pearson correlation of 0.881 between irradiance and power.

The Optuna-optimized MLP neural network with three hidden layers of 256
units, a dropout probability of 0.355, and a learning rate of 0.00579
achieved an R² of 0.954, RMSE of 5.17 kW, and MAE of 3.72 kW on the
held-out test set. This represents a 13.91 percentage point improvement
in R² and a 51.4\% reduction in RMSE compared to the best linear
regression baseline (R² = 0.815, RMSE = 10.65 kW). The mean residual of
0.281 kW confirms approximately unbiased predictions. These results
demonstrate that an ANOVA-guided, hyperparameter-optimized neural
network can deliver highly accurate hourly PV forecasts from minimal
input features, achieving performance competitive with or superior to
more complex deep learning architectures reported in recent literature.

Several limitations of this study warrant acknowledgment. The dataset
originates from a single PV plant, and the generalizability of the
optimized hyperparameters and model architecture to other plants with
different capacities, tilt angles, or climatic conditions requires
validation through transfer learning experiments. The input feature set
is limited to POA irradiance and temporal features; incorporating
temperature, humidity, and cloud cover data could further improve
performance, particularly during periods of complex weather dynamics.
The MLP architecture does not explicitly model temporal dependencies
between consecutive hours, and LSTM-based sequential forecasting may
offer additional improvement for multi-step-ahead predictions.

Future work will extend this framework in several directions. First,
ensemble methods including Random Forest, XGBoost, and stacking
approaches will be benchmarked against the MLP to provide a
comprehensive model comparison across algorithmic families. Second, LSTM
and GRU recurrent architectures will be evaluated for multi-step
sequential forecasting. Third, transfer learning experiments across
multiple PV plant locations will assess the portability of trained
models. Fourth, real-time deployment considerations including model
update frequency, computational latency, and integration with SCADA
systems will be investigated to bridge the gap between research accuracy
and operational utility.

\begin{center}\rule{0.5\linewidth}{0.5pt}\end{center}

\subsection{Acknowledgements}\label{acknowledgements}

The authors gratefully acknowledge the data provided by the operational
solar PV plant facility. This research was conducted using open-source
software tools including Python 3.x, PyTorch for neural network
implementation, Optuna for hyperparameter optimization, scikit-learn for
linear regression modeling, pandas for data manipulation, and Matplotlib
and Seaborn for visualization. The computational work was performed on
standard workstation hardware without requiring GPU acceleration.

\begin{center}\rule{0.5\linewidth}{0.5pt}\end{center}

\subsection{References}\label{references}

{[}1{]} M. Elsaraiti and A. Merabet, “Solar power forecasting using deep
learning techniques,” \emph{IEEE Access}, vol.~10, pp.~31692-31698,
2022. doi: 10.1109/ACCESS.2022.3160484

{[}2{]} International Energy Agency, “World Energy Outlook 2023,” IEA,
Paris, 2023. {[}Online{]}. Available:
https://www.iea.org/reports/world-energy-outlook-2023

{[}3{]} R. Asghar, F. R. Fulginei, M. Quercio, and A. Mahrouch,
“Artificial neural networks for photovoltaic power forecasting: A review
of five promising models,” \emph{IEEE Access}, vol.~12, pp.~89907-89932,
2024. doi: 10.1109/ACCESS.2024.3420693

{[}4{]} M. Husein, E. J. Gago, B. Hasan, and M. C. Pegalajar, “Towards
energy efficiency: A comprehensive review of deep learning-based
photovoltaic power forecasting strategies,” \emph{Heliyon}, vol.~10, no.
13, e34419, 2024. doi: 10.1016/j.heliyon.2024.e33419

{[}5{]} S. Al-Dahidi, M. Madhiarasan, and L. Al-Ghussain, “Forecasting
solar photovoltaic power production: A comprehensive review and
innovative data-driven modeling framework,” \emph{Energies}, vol.~17,
no. 16, p.~4145, 2024. doi: 10.3390/en17164145

{[}6{]} W. C. Tsai, C. S. Tu, C. M. Hong, and W. M. Lin, “A review of
state-of-the-art and short-term forecasting models for solar PV power
generation,” \emph{Energies}, vol.~16, no. 14, p.~5436, 2023. doi:
10.3390/en16145436

{[}7{]} V. Prema, M. S. Bhaskar, D. Almakhles, N. Gowtham et al.,
“Critical review of data, models and performance metrics for wind and
solar power forecast,” \emph{IEEE Access}, vol.~10, pp.~667-688, 2021.
doi: 10.1109/ACCESS.2021.3137419

{[}8{]} F. Pandžić and T. Capuder, “Advances in short-term solar
forecasting: A review and benchmark of machine learning methods and
relevant data sources,” \emph{Energies}, vol.~17, no. 1, p.~97, 2024.
doi: 10.3390/en17010097

{[}9{]} G. Sahin, G. Isik, and W. G. van Sark, “Predictive modeling of
PV solar power plant efficiency considering weather conditions: A
comparative analysis of artificial neural networks and multiple linear
regression,” \emph{Energy Reports}, vol.~10, pp.~2930-2946, 2023. doi:
10.1016/j.egyr.2023.09.097

{[}10{]} A. Mellit, A. M. Pavan, and V. Lughi, “Deep learning neural
networks for short-term photovoltaic power forecasting,” \emph{Renewable
Energy}, vol.~178, pp.~228-235, 2021. doi: 10.1016/j.renene.2021.02.166

{[}11{]} M. N. Akhter, S. Mekhilef, H. Mokhlis, and Z. M. Almohaimeed,
“An hour-ahead PV power forecasting method based on an RNN-LSTM model
for three different PV plants,” \emph{Energies}, vol.~15, no. 6,
p.~2243, 2022. doi: 10.3390/en15062243

{[}12{]} I. K. Bazionis, M. A. Kousounadis-Knousen et al., “A taxonomy
of short-term solar power forecasting: Classifications focused on
climatic conditions and input data,” \emph{IET Renewable Power
Generation}, vol.~17, no. 10, pp.~2539-2556, 2023. doi:
10.1049/rpg2.12736

{[}13{]} D. A. Wood, “Hourly-averaged solar plus wind power generation
for Germany 2016: Long-term prediction, short-term forecasting, data
mining and outlier analysis,” \emph{Sustainable Cities and Society},
vol.~60, p.~102230, 2020. doi: 10.1016/j.scs.2020.102227

{[}14{]} M. S. Hossain and H. Mahmood, “Short-term photovoltaic power
forecasting using an LSTM neural network and synthetic weather
forecast,” \emph{IEEE Access}, vol.~8, pp.~172524-172533, 2020. doi:
10.1109/ACCESS.2020.3024901

{[}15{]} C. H. Liu, J. C. Gu, and M. T. Yang, “A simplified LSTM neural
networks for one day-ahead solar power forecasting,” \emph{IEEE Access},
vol.~9, pp.~17174-17195, 2021. doi: 10.1109/ACCESS.2021.3053638

{[}16{]} X. Li, L. Ma, P. Chen, H. Xu, Q. Xing, J. Yan, S. Lu, and H.
Fan, “Probabilistic solar irradiance forecasting based on XGBoost,”
\emph{Energy Reports}, vol.~8, pp.~1087-1095, 2022. doi:
10.1016/j.egyr.2022.02.251

{[}17{]} D. Salman, C. Direkoglu, and M. Kusaf, “Hybrid deep learning
models for time series forecasting of solar power,” \emph{Neural
Computing and Applications}, vol.~36, pp.~9191-9213, 2024. doi:
10.1007/s00521-024-09558-5

{[}18{]} M. Al Hazza, H. Attia, and K. Hossin, “Solar photovoltaic power
prediction using a statistical approach based on analysis of variance,”
\emph{Journal of Solar Energy and Sustainable Development}, vol.~13, no.
1, pp.~37-50, 2024.

{[}19{]} V. Demir and H. Citakoglu, “Forecasting of solar radiation
using different machine learning approaches,” \emph{Neural Computing and
Applications}, vol.~35, pp.~7329-7350, 2023. doi:
10.1007/s00521-022-07841-x

{[}20{]} C. Van Zyl, X. Ye, and R. Naidoo, “Harnessing eXplainable
artificial intelligence for feature selection in time series energy
forecasting: A comparative analysis of Grad-CAM and SHAP,” \emph{Applied
Energy}, vol.~355, p.~122314, 2024. doi: 10.1016/j.apenergy.2023.122079

{[}21{]} D. K. Dhaked, S. Dadhich, and D. Birla, “Power output
forecasting of solar photovoltaic plant using LSTM,” \emph{Green Energy
and Intelligent Transportation}, vol.~2, no. 5, p.~100113, 2023. doi:
10.1016/j.geits.2023.100113

{[}22{]} P. Li, K. Zhou, X. Lu, and S. Yang, “A hybrid deep learning
model for short-term PV power forecasting,” \emph{Applied Energy},
vol.~259, p.~114216, 2020. doi: 10.1016/j.apenergy.2019.114216

{[}23{]} M. N. Akhter, S. Mekhilef, H. Mokhlis, R. Ali, and M. Usama, “A
hybrid deep learning method for an hour ahead power output forecasting
of three different photovoltaic systems,” \emph{Applied Energy},
vol.~307, p.~118185, 2022. doi: 10.1016/j.apenergy.2021.118185

{[}24{]} M. Konstantinou, S. Peratikou, and A. G. Charalambides, “Solar
photovoltaic forecasting of power output using LSTM networks,”
\emph{Atmosphere}, vol.~12, no. 1, p.~124, 2021. doi:
10.3390/atmos12010124

\end{document}
